\chapter{事件、国家、主题与图件索引}
\label{app:index}

本索引选择回到首次深描的事件节点。每一条路由都由事件通向政权、主题、来源或图件，并可从相应附录再回到叙事；它不是一份逐条年表。

\section{事件索引}
\label{app:index:events}

\begin{description}
  \item[魏] \eventindexentry{TK-0220-WEI-LUOYANG}{魏定都洛阳}；\eventindexentry{TK-0238-SIMA-YI-LIAODONG}{司马懿平辽东}；\eventindexentry{TK-0249-GAOPING-TOMBS}{高平陵政变}；\eventindexentry{TK-0258-ZHUGE-DAN-DEFEAT}{寿春陷落}；\eventindexentry{TK-0260-CAO-MAO-KILLED}{曹髦之死}；\eventindexentry{TK-0265-WEI-ABDICATION}{魏晋禅代}。
  \item[蜀汉] \eventindexentry{TK-0221-SHU-HAN-FOUNDATION}{刘备在成都称帝}；\eventindexentry{TK-0223-SHU-WU-RECONCILIATION}{蜀吴复通}；\eventindexentry{TK-0225-SHU-NANZHONG}{南征}；\eventindexentry{TK-0227-ZHUGE-LIANG-HANZHONG}{组织北伐}；\eventindexentry{TK-0228-JIETING}{街亭失守}；\eventindexentry{TK-0234-WUZHANG}{五丈原撤军}；\eventindexentry{TK-0263-CHENGDU-SURRENDER}{刘禅出降}。
  \item[孙吴] \eventindexentry{TK-0222-WEI-INVESTS-SUN-QUAN}{魏册孙权}；\eventindexentry{TK-0229-SUN-QUAN-EMPEROR}{孙权称帝}；\eventindexentry{TK-0229-JIANYE-CAPITAL}{迁都建业}；\eventindexentry{TK-0252-DONGXING}{东兴之战}；\eventindexentry{TK-0272-XILING}{西陵之战}。
  \item[跨政权与接收] \eventindexentry{TK-0222-YILING}{夷陵之战}；\eventindexentry{TK-0238-HIMIKO-EMBASSY}{魏受倭使}；\eventindexentry{TK-0264-ZHONG-HUI-REBELLION}{成都兵变}；\eventindexentry{TK-0280-SUN-HAO-SURRENDER}{孙皓出降}。
  \item[书写与法律] \eventindexentry{TK-0240-ZHENGSHI-STONE-CLASSICS}{正始三体石经}；\eventindexentry{TK-0268-TAISHI-CODE}{《泰始律》的颁行}。
\end{description}

\section{国家与主题索引}
\label{app:index:themes}

\begin{description}
  \item[魏：洛阳、辅政与淮南] \eventref{TK-0220-WEI-LUOYANG}{魏定都洛阳}、\eventref{TK-0249-GAOPING-TOMBS}{高平陵政变}与\eventref{TK-0265-WEI-ABDICATION}{魏晋禅代}。
  \item[蜀汉：成都、南中与汉中] \eventref{TK-0221-SHU-HAN-FOUNDATION}{蜀汉建国}、\eventref{TK-0225-SHU-NANZHONG}{南征}、\eventref{TK-0228-JIETING}{街亭失守}与\eventref{TK-0263-CHENGDU-SURRENDER}{成都出降}。
  \item[孙吴：长江、继承与峡口] \eventref{TK-0229-SUN-QUAN-EMPEROR}{孙权称帝}、\eventref{TK-0252-DONGXING}{东兴}与\eventref{TK-0272-XILING}{西陵}。
  \item[战争道路与边缘网络] \eventref{TK-0222-YILING}{夷陵}、\eventref{TK-0238-HIMIKO-EMBASSY}{倭使}与\eventref{TK-0280-SUN-HAO-SURRENDER}{灭吴}。
  \item[文本、法律与地方尺度] \eventref{TK-0240-ZHENGSHI-STONE-CLASSICS}{正始三体石经}与\eventref{TK-0268-TAISHI-CODE}{《泰始律》}。
\end{description}

\section{来源索引}
\label{app:index:sources}

\begin{description}
  \sourceindexentry{\hyperref[src:chronicles]{三国志、裴松之注与资治通鉴}}{共同年序与并行政治的入口；回到\eventref{TK-0222-YILING}{夷陵}、\eventref{TK-0249-GAOPING-TOMBS}{高平陵}和\eventref{TK-0263-CHENGDU-SURRENDER}{成都出降}。}
  \sourceindexentry{\hyperref[src:chronicles]{后汉书与晋书}}{禅代、受降和统一的路径；回到\eventref{TK-0265-WEI-ABDICATION}{魏晋禅代}与\eventref{TK-0280-SUN-HAO-SURRENDER}{孙皓出降}。}
  \sourceindexentry{\hyperref[src:documents]{走马楼吴简、石经与律文}}{地方文书、官学与规则的不同尺度；回到\eventref{TK-0240-ZHENGSHI-STONE-CLASSICS}{石经}与\eventref{TK-0268-TAISHI-CODE}{律文}。}
  \sourceindexentry{\hyperref[src:regional]{华阳国志、区域史与考古}}{道路与边缘的比较路径；回到\eventref{TK-0225-SHU-NANZHONG}{南征}、\eventref{TK-0238-SIMA-YI-LIAODONG}{辽东}与\eventref{TK-0238-HIMIKO-EMBASSY}{倭使}。}
\end{description}

\section{图件索引}
\label{app:index:figures}

\begin{description}
  \figureindexentry{fig:vol02b:period-timeline}{三国至西晋统一的共同时间轴（220--280年）}
  \figureindexentry{fig:vol02b:polity-corridors}{三国政治中心与主要竞争走廊}
  \figureindexentry{fig:vol02b:yiling-river-defense}{夷陵、白帝与孙吴上游防务；见\eventref{TK-0222-YILING}{夷陵之战}}
  \figureindexentry{fig:vol02b:northern-liaodong}{陇右、辽东与淮南的边缘压力；见\eventref{TK-0238-SIMA-YI-LIAODONG}{平辽东}与\eventref{TK-0238-HIMIKO-EMBASSY}{倭使}}
  \figureindexentry{fig:vol02b:jin-conquest-wu}{西晋伐吴的多路推进；见\eventref{TK-0279-JIN-INVADES-WU}{多路伐吴}与\eventref{TK-0280-SUN-HAO-SURRENDER}{孙皓出降}}
\end{description}
